\documentclass[11pt,a4paper]{article}

%% PACHETE MATEMATICE
\usepackage{amsmath,amssymb,amsfonts}
\usepackage{amsthm}
\usepackage{mathpazo}
\usepackage{eulervm}
\usepackage{enumerate}
\usepackage{color}
\usepackage{graphicx}
\usepackage[all]{xy}
\usepackage{xcolor}
\usepackage{framed}
\usepackage[unicode]{hyperref}
\setcounter{MaxMatrixCols}{20}
\usepackage{multicol}
%\usepackage[romanian]{babel}


%% ASPECT
\linespread{1.05}
\usepackage[T1]{fontenc}
\usepackage[utf8x]{inputenc}
\usepackage[margin=2cm]{geometry}
% \usepackage[color]{showkeys}
\usepackage{anyfontsize}
\usepackage{titlesec}
\titleformat*{\section}{\large\bfseries}

\usepackage{fancyhdr}
\pagestyle{fancy}
\rhead{\small \color{gray} Notițe de Adrian Manea}
\lhead{\small \color{gray} ETTI-AG, 2017---2018, A. Niță \& A. Oprina}


\usepackage[autostyle = false, style = german]{csquotes}
\usepackage[romanian]{babel}
\newcommand\qq{\enquote}

%% COMENZILE MELE
\renewcommand*{\proofname}{Dem.:}
\newcommand\bb{\mathbb}
\newcommand\dr{\mathrm}
\newcommand\kal{\mathcal}
\newcommand\fk{\mathfrak}
\newcommand\op{\oplus}
\newcommand\ot{\otimes}
\newcommand\ds{\displaystyle}
\newcommand\id{\indent}
\newcommand\nid{\noindent}
\newcommand\seq{\subseteq}
\newcommand\sneq{\subsetneq}
\newcommand\speq{\supseteq}
\newcommand\spneq{\supsetneq}
\newcommand\rar{\rightarrow}
\newcommand\Rar{\Rightarrow}
\newcommand\xrar{\xrightarrow}
\newcommand\lar{\leftarrow}
\newcommand\Lar{\Leftarrow}
\newcommand\xlar{\xleftarrow}
\newcommand\lrar{\leftrightarrow}
\newcommand\Lrar{\Leftrightarrow}
\newcommand\ideal{\trianglelefteq}
\newcommand\ai{\text{ a.\^{i}. }}
\newcommand\sk{(\textit{Schi\c{t}\u{a}}) }
\newcommand\rhup{\rightharpoonup}
\newcommand\rhdn{\rightharpoondown}
\newcommand\lhup{\leftharpoonup}
\newcommand\lhdn{\leftharpoondown}
\newcommand\vid{\emptyset}
\newcommand\wh{\widehat}
\newcommand\ol{\overline}
\newcommand\NN{\mathbb{N}}
\newcommand\RR{\mathbb{R}}
\newcommand\ZZ{\mathbb{Z}}
\newcommand\QQ{\mathbb{Q}}
\newcommand\CC{\mathbb{C}}

%% STILURILE MELE
\newtheoremstyle{dotless}{}{}{\itshape}{}{\bfseries}{:}{ }{}

\theoremstyle{dotless}
\newtheorem{theorem}{Teorem\u{a}}[section]
\newtheorem{proposition}{Propozi\c{t}ie}[section]
\newtheorem{lemma}{Lem\u{a}}[section]
\newtheorem{corollary}{Corolar}[section]
\newtheorem{exercise}{Exerci\c{t}iu}

\newtheoremstyle{dotlessdr}{}{}{}{}{\bfseries}{:}{ }{}
\theoremstyle{dotlessdr}
\newtheorem{example}{Exemplu}[section]
\newtheorem{definition}{Defini\c{t}ie}[section]
\newtheorem{remark}{Observa\c{t}ie}[section]




\begin{document}

\begin{center}
  {\large\textbf{Seminar 6 \\
    Spații euclidiene}}
\end{center}

\vspace{1cm}

\section{Produs scalar}
\label{sec:teorie}

Spațiile euclidiene lărgesc definițiile spațiilor vectoriale generale,
adăugînd construcții foarte utile în perspectiva aplicațiilor în geometrie.

Prima definiție importantă este:

\begin{definition}\label{def:prod-sc}
  Fie $V$ un spațiu vectorial real. Se numește \emph{produs scalar} pe $V$
  o aplicație $\langle \cdot, cdot \rangle : V \times V \to \RR$, care, pentru orice
  vectori $x, y, z \in V$ și orice scalar $\alpha \in \RR$ are proprietățile:
  \begin{itemize}
  \item \emph{comutativitate:} $\langle x, y \rangle = \langle y, x \rangle$;
  \item \emph{aditivitate:} $\langle x + z, y \rangle = \langle x, y \rangle + %
    \langle z, y \rangle$;
  \item \emph{balansare:} $\langle \alpha x, y \rangle = \alpha \langle x, y \rangle$;
  \item \emph{pozitivitate:} $\langle x, x \rangle = 0 \Leftrightarrow x = 0$.
  \end{itemize}
\end{definition}

Cu această construcție, avem:

\begin{definition}\label{def:preh-euc}
  Un spațiu vectorial real pe care s-a definit un produs scalar se numește
  \emph{spațiu prehilbertian (real)}. \footnote{Denumirea provine de la %
    matematicianul german David Hilbert (1862---1943), iar particular \emph{pre-} arată %
    că aceste spații sînt utilizate pentru a obține \emph{spațiile Hilbert}, extrem %
    de importante în fizica teoretică. Ele adaugă și proprietăți de analiză %
    vectorilor, studiind, de fapt, \emph{spații vectoriale de funcții}.}

  Dacă spațiul este finit dimensional, el se numește \emph{euclidian}.
\end{definition}

În continuare, vom presupune că spațiul $V$ pe care îl folosim este un spațiu
euclidian.

Ca în geometrie, numim doi vectori $x, y \in V$ \emph{ortogonali} dacă
$\langle x, y \rangle = 0$. Vom mai nota aceasta și prin $x \perp y$.

Folosind noțiunea de produs scalar, putem defini, pentru orice vectori $x \in V$,
aplicația \emph{normă}, definită prin:
\[
  || x || = \sqrt{ \langle x, x \rangle }.
\]
Această aplicație generalizează noțiunea de \emph{modul} al unui vector.

Exemplele pe care le vom întîlni cel mai des sînt:
\begin{itemize}
\item Pentru spațiul euclidian $\RR^n$, produsul scalar standard este:
  \[
    \langle x, y \rangle = \sum_{i = 1}^n x_i \cdot y_i,
  \]
  unde $x = (x_1, \dots, x_n)$, iar $y = (y_1, \dots, y_n)$;
\item Pentru spațiul de matrice $M_{m,n}(\RR)$, putem defini produsul scalar:
  \[
    \langle A, B \rangle = \mathrm{tr}(A^t \cdot B).
  \]
\item Pentru $V = C^0([a, b])$, spațiul vectorial al tuturor funcțiilor continue
  de forma $f : [a, b] \to \RR$, definim:
  \[
    \langle f, g \rangle = \int_a^b f(x) \cdot g(x) dx,
  \]
  structură cu care obținem un spațiu prehilbertian infinit dimensional.
\end{itemize}

În contextul unui spațiu euclidian, au loc relațiile cunoscute din cazul
geometriei plane:
\begin{itemize}
\item $\forall x,y \in V, |\langle x, y \rangle| \leq ||x|| \cdot ||y||$ (inegalitatea
  lui Cauchy-Buniakowski-Schwartz);
\item $\forall x, y \in V, ||x + y||^2 + ||x - y||^2 = 2(||x||^2 + ||y||^2)$ (regula %
  paralelogramului);
\item Dacă $x \perp y$, atunci $||x + y||^2 = ||x||^2 + ||y||^2$ (teorema lui Pitagora).
\end{itemize}

În cazul spațiilor complexe, se impun modificări la definiții.

\begin{definition}\label{def:prod-sc-compl}
  Fie $V$ un spațiu vectorial complex. Se numește \emph{produs scalar} pe $V$
  o aplicație $\langle \cdot, \cdot \rangle : V \times V \to \CC$,
  care, pentru orice vectori $x, y, z \in V$ și orice scalar $\alpha \in \CC$,
  are proprietățile:
  \begin{itemize}
  \item $\langle x, y \rangle = \overline{\langle y, x \rangle}$, unde prin $\overline{()}$
    am notat conjugata complexă;
  \item $\langle x + z, y \rangle = \langle x, y \rangle + \langle z, y \rangle$;
  \item $\langle \alpha x, y \rangle = \alpha \langle x, y \rangle$;
  \item $\langle x, x \rangle \geq 0$, cu egalitate dacă și numai dacă $x = 0$.
  \end{itemize}
\end{definition}

\begin{remark}\label{rk:conj}
  Remarcați că, folosind prima și a treia proprietate, obținem:
  \[
    \langle x, \alpha y \rangle = \overline{\alpha} \langle x, y \rangle.
  \]
\end{remark}

Cu aceasta, avem:

\begin{definition}\label{def:sp-unit}
  Un spațiu vectorial complex, pe care s-a definit un produs scalar se numește
  \emph{spațiu prehilbertian complex}.

  Spațiile prehilbertiene complexe cu dimensiune finită se numesc \emph{spații unitare}.
\end{definition}

Exemplele din cazul real pot fi adaptate:
\begin{itemize}
\item Fie $V = \CC^n$. Atunci, cu notațiile de mai sus, produsul scalar se
  poate defini prin:
  \[
    \langle x, y \rangle = \sum_{i = 1}^n x_i \cdot \overline{y_i}.
  \]
\item În spațiul vectorial complex al funcțiilor continue $f : [a, b] \to \CC$,
  putem defini produsul scalar complex prin:
  \[
    \langle f, g \rangle = \int_a^b f(x) \overline{g(x)} dx.
  \]
\end{itemize}

%%%%%%%%%%%%%%%%%%%%%%%%%%%%%%%%%%%%%%%%%%%%%%%%%%%%%%%%%%%%%%%%%%%%%%

\section{Ortonormare}
\label{sec:orton}

Date aceste noțiuni, pornind cu conceptul de produs scalar, utilizările
pe care le vom studia sînt în următoarea direcție:

\begin{definition}\label{def:orton}
  Fie $V$ un spațiu prehilbertian (real sau complex).
  \begin{itemize}
  \item Un sistem finit sau numărabil $\{ e_i \}$ de vectori din $V$ se
    numește \emph{sistem ortogonal} dacă toți vectorii sînt nenuli și
    ortogonali doi cîte doi.
  \item Sistemul se numește \emph{ortonormat} dacă este ortogonal, iar
    norma (\qq{lungimea}) fiecărui vector este 1.
  \end{itemize}
\end{definition}

Cazul care ne interesează este acela în care sistemul de vectori de mai
sus este o \emph{bază} pentru $V$, situație în care se numește
\emph{bază ortonormată}. \footnote{Acesta este, de fapt, cazul din geometria
  tridimensională, \qq{clasică}. Baza este dată de versorii $\vec{i}, \vec{j}, %
  \vec{k}$, care sînt ortonormați.}

Dată o bază \qq{obișnuită} a unui spațiu vectorial (euclidian sau unitar), din
ea se poate obține una \emph{ortonormată}, folosind \textbf{procedeul Gram-Schmidt}.

Fie $B = \{ u_1, \dots, u_n \}$ o bază oarecare în $V$. Construim din ea o bază
\emph{ortogonală} $B_1 = \{ v_1, \dots, v_n \}$, apoi o normăm, înmulțind fiecare
vector cu inversul normei sale, obținînd baza \emph{ortonormată} %
$B_2 = \{ \frac{v_i}{||v_i||} \}$.

Procedeul pornește astfel:
\begin{itemize}
\item Luăm $v_1 = u_1$;
\item Mai departe, definim $v_2 = u_2 + \alpha v_1$, cu $\alpha$ ales astfel
  încît $v_2 \perp v_1$. Din această condiție, găsim:
  \[
    \alpha = - \frac{\langle u_2, v_1 \rangle}{\langle v_1, v_1 \rangle}.
  \]
\item În continuare, luăm $v_3 = u_3 + \alpha_1 v_1 + \alpha_2 v_2$, cu $\alpha_{1,2}$
  astfel încît $v_3 \perp v_1$ și $v_3 \perp v_2$. Se obțin:
  \[
    \alpha_1 = -\frac{\langle u_3, v_1 \rangle}{\langle v_1, v_1 \rangle}, \quad %
    \alpha_2 = -\frac{\langle u_3, v_2 \rangle}{\langle v_2, v_2 \rangle}.
  \]
\item Procedeul continuă similar pînă obținem baza $B_1$, pe care apoi o normăm.
\end{itemize}

%%%%%%%%%%%%%%%%%%%%%%%%%%%%%%%%%%%%%%%%%%%%%%%%%%%%%%%%%%%%%%%%%%%%%%

\section{Complement ortogonal}
\label{sec:compl-ort}

Extindem noțiunea de ortogonalitate la spații.

\begin{definition}\label{def:compl-ort}
  Fie $E$ un spațiu euclidian (sau unitar) cu $\dim E = n < \infty$ și $V$ un
  subspațiu vectorial al lui $E$, cu $\dim V = p$, astfel încît $1 \leq p \leq n - 1$.

  Se numește \emph{complementul ortogonal} al lui $V$, notat $V^\perp$, un
  subspațiu $V^\perp \seq E$, astfel încît $V \oplus V^\perp = E$ și
  $V \perp V^\perp$, adică $\forall x \in V, \forall y \in V^\perp, x \perp y$.
\end{definition}

În calcule, ne vom baza pe următoarea:

\begin{theorem}\label{thm:comp-ort}
  Fie $E$ un spațiu euclidian (unitar). Pentru orice subspațiu vectorial nenul
  $V \seq E$, există și este unic complementul ortogonal $V^\perp$ al lui $V$.
\end{theorem}

Așadar, următoarele proprietăți vor fi fundamentale în exercițiile de
determinare a complementului ortogonal al unui spațiu. Fie $S \hookrightarrow V$.
Atunci:
\begin{itemize}
\item $S \oplus S^\perp \simeq V$, deci (cf. Grassmann):
  \[
    \dim S + \dim S^\perp = \dim V.
  \]
\item $(S^\perp)^\perp \simeq S$;
\item Putem scrie matricea formată din vectorii bazei lui $S$ și căutăm vectorii
  bazei lui $S^\perp$ (tot matriceal), astfel încît produsul celor două matrice
  să fie nul. Problema se reduce, astfel, la un sistem de ecuații omogene.
\end{itemize}

Avem nevoie și de:

\begin{definition}\label{def:end-ort}
  Fie $E$ un spațiu euclidian (unitar). Un endomorfism $f : E \to E$ se numește
  \emph{ortogonal (unitar)} dacă duce baze ortonormate în baze ortonormate.
\end{definition}

Pentru ele, avem următoarea:

\begin{theorem}\label{thm:car-end-ort}
  Dacă $E$ este un spațiu euclidian (unitar), iar $f : E \to E$ este un endomorfism,
  următoarele afirmații sînt echivalente:
  \begin{enumerate}[(a)]
  \item $f$ este ortogonal (unitar);
  \item $\forall x, y \in E, \langle x, y \rangle = \langle f(x), f(y) \rangle$;
  \item matricea lui $f$, scrisă în orice bază ortonormată, este inversabilă,
    iar $(M_f^B)^{-1} = (M_f^B)^t$ (respectiv $\overline{(M_f^B)^t}$ pentru cazul
    complex).
  \end{enumerate}
\end{theorem}

De fapt, putem defini în general:

\begin{definition}\label{def:mat-sim-ort}
  O matrice $A \in M_n(\RR)$ se numește \emph{simetrică} atunci cînd $A = A^t$
  și \emph{ortogonală} cînd $A^{-1} = A^t$.

  Pentru cazul complex, o matrice $A \in M_n(\CC)$ se numește \emph{hermitică}
  \footnote{Charles Hermite (1822--1901) %
    {\color{blue} \href{https://en.wikipedia.org/wiki/Charles_Hermite}{[wiki]}}.}
  atunci cînd $A = \overline{A^t}$ și \emph{unitară} dacă este inversabilă
  și $A^{-1} = \overline{A^t}$.
\end{definition}

Se poate demonstra simplu că:
\begin{itemize}
\item Matricea asociată unui endomorfism ortogonal într-o bază ortonormată
  este ortogonală;
\item Dacă $A$ este ortogonală, atunci $\det A = \pm 1$;
\item Matricea asociată unui endomorfism unitar într-o bază ortonormată este
  unitară;
\item Dacă $A$ este unitară, atunci $|\det A| = 1$.
\end{itemize}

%%%%%%%%%%%%%%%%%%%%%%%%%%%%%%%%%%%%%%%%%%%%%%%%%%%%%%%%%%%%%%%%%%%%%%

\section{Exerciții}
\label{sec:exc}

\indent\indent 1. Aplicați algoritmul Gram-Schmidt pentru:
\begin{enumerate}[(a)]
\item Baza $\kal{B} = \{(1, 1, -1), (3, -1, -1), (0, -1, 1) \}$ a spațiului $\RR^3$;
\item Baza $\kal{B} = \{(0, 1, -1), (2, 0, 1), (-1, 1, -1) \}$ a spațiului $\RR^3$.
\end{enumerate}

\vspace{1cm}

2. Fie $V = \RR_3[X]$. Definim aplicația $\langle \cdot, \cdot \rangle : V \times V \to \RR$
prin:
\[
  \langle p, q \rangle = \sum_{k = 0}^3 (k!)^2 a_k b_k,
\]
unde $p = \sum a_i X^i$ și $q = \sum b_j X^j$.

\begin{enumerate}[(a)]
\item Arătați că aplicația este un produs scalar;
\item Calculați $||h||$, unde $h = 1 + 5x - 4x^2 + 5x^3$.
\end{enumerate}

\vspace{1cm}

3. Fie $C = \kal{C}^\infty([1, e])$ și aplicația definită prin:
\[
  \langle f, g \rangle = \int_1^e f(x) g(x) \ln x dx.
\]
\begin{enumerate}[(a)]
\item Arătați că aplicația este un produs scalar;
\item Calculați $||h||$, pentru $h : [1, e] \to \RR, h(x) = \sqrt{x}$;
\item Aflați o funcție de gradul întîi $g$ care să fie perpendiculară pe
  funcția constantă $f(x) = 5$ (în raport cu produsul scalar de mai sus).
\end{enumerate}

\vspace{1cm}

4. Fie $S = \{ (x, y, z) \in \RR^3 \mid x + y = 0 \}$.
\begin{enumerate}[(a)]
\item Determinați $S^\perp$;
\item Fie vectorul $v = (1, 4, 5)$. Găsiți coordonatele lui $v$ în $S$,
  respectiv în $S^\perp$.
\end{enumerate}

\vspace{1cm}

5. Calculați complementul ortogonal al subspațiului al $\RR^4$, generat
peste $\RR$ de vectorii
\[
  (1, 3, 0, 2), (3, 7, -1, 2), (2, 4, -1, 0).
\]

\vspace{1cm}

6. În spațiul $\RR_2[X]$ se consideră vectorii:
\begin{align*}
  p_1 &= 3x^2 + 2x + 1 \\
  p_2 &= 3x^2 + 2x + 5 \\
  p_3 &= -x^2 + 2x + 1 \\
  p_4 &= 3x^2 + 5x + 2.
\end{align*}

Să se determine un polinom $p$ echidistant de vectorii $p_1, p_2, p_3, p_4$,
în raport cu distanța euclidiană.

\vspace{1cm}

7. În spațiul $\RR^4$, se consideră vectorii:
\[
  x = (1, 0, 1, 3), \quad y = (-1, 1, 1, 0).
\]
Să se completeze acești vectori pînă la o bază ortogonală, apoi să se normeze baza
rezultată.

\vspace{1cm}

8. Se consideră subspațiul $S = \mathrm{Sp}(x_1, x_2, x_3) \hookrightarrow \RR^4$,
unde:
\[
  x_1 = (1, 1, 2, 1), x_2 = (-1, 0, 2, 3), x_3 = (1, 2, -1, -3).
\]

\begin{enumerate}[(a)]
  \item Găsiți un vector $v \in \RR^4$, astfel încît $v \perp S$;
  \item Găsiți $S^\perp$.
\end{enumerate}
  
\vspace{1cm}

9. Să se determine o bază ortonormată a subspațiului:
\[
  S = \{(x_1, x_2, x_3, x_4) \in \RR^4 \mid x_1 - x_2 + 2x_3 - x_4 = 0 \} %
  \hookrightarrow \RR^4.
\]

\end{document}
